% Chapter: The Cohomological Calculus of Constructions (C⁴)
% Rewritten to fit the Synergy Atlas — refinement site, descent=HComp,
% two dimension systems, cover algebra.

\chapter{The Cohomological Calculus of Constructions}
\label{ch:c4-calculus}

$\mathrm{C}^4$ is a dependent type theory built on three pillars:
the Calculus of Inductive Constructions (CIC), cubical type theory,
and sheaf cohomology over a \emph{refinement site}.  The central
discovery is that the cubical and cohomological structures are not
merely co-present but \emph{unified}: the descent rule that glues
sheaf sections is the same operation as the Kan filling rule that
fills cubical horns.  This chapter develops that unification from first
principles, starting with the two independent dimension systems that
together give $\mathrm{C}^4$ its distinctive power.


% ─────────────────────────────────────────────────────────
\section{The Fundamental Identification: Descent \(\equiv\) HComp}
\label{sec:descent-hcomp-id}

Before defining any syntax, we state the organizing principle.

\begin{theorem}[Fundamental Identification]
\label{thm:fundamental-id}
Homogeneous composition (\(\hfilling\)) in the cubical fragment and
cohomological descent in the sheaf fragment are two manifestations of
a single universal construction: the \textbf{filling of a partial cube
from compatible face data}.
\end{theorem}

The identification is made precise by a common abstraction:

\begin{definition}[Partial Boundary Data]
A \textbf{partial boundary} for a type \(A\) on a cover
\(\covers = \{U_\alpha\}_{\alpha \in I}\) is a family of local sections
\(s_\alpha \in A(U_\alpha)\) together with compatibility data:
\[
  \forall \alpha, \beta.\;
  \restr{s_\alpha}{U_\alpha \cap U_\beta}
  \;=\;
  \restr{s_\beta}{U_\alpha \cap U_\beta}.
\]
\end{definition}

In the \textbf{cubical reading}, the \(U_\alpha\) are the faces
\([i_k = \varepsilon]\) of the interval \(\interval^n\), the local
sections are the face terms, and the compatibility data is edge
agreement.  The filled section is the term inhabiting the interior.

In the \textbf{cohomological reading}, the \(U_\alpha\) are the open
sets of the refinement cover \(\covers\), the local sections are
per-fiber proofs, and the compatibility data is the overlap proofs.
The filled section is the global proof.

\begin{proof}[Proof of Theorem~\ref{thm:fundamental-id}]
We exhibit a single type-theoretic rule that covers both cases.
Let \(\mathcal{F}\) be a sheaf (in the cubical case, the constant
sheaf \(\underline{A}\); in the cohomological case, a sheaf of
refinement types).

The \textbf{filling condition} is: given partial boundary data as
above, there exists a unique global section
\(s \in \mathcal{F}(X)\) with \(\restr{s}{U_\alpha} = s_\alpha\).

The existence of \(s\) is equivalent to
\(\check{H}^1(\covers, \underline{\mathrm{Aut}}(\mathcal{F})) = 0\)
(SGA~1).

For our two settings:
\begin{enumerate}
\item \textbf{Cubical (\(\hfilling\)):}  The Kan condition asserts
  that every horn has a filler, which is exactly
  \(\check{H}^1(\partial\interval^n, \underline{A}) = 0\) for
  contractible \(\partial\interval^n\).

\item \textbf{Cohomological (descent):}  The propositional coefficients
  of our sheaves have \(\check{H}^1 = 0\) by Hedberg's theorem
  (Theorem~\ref{thm:hedberg-collapse} below), so every compatible
  boundary data has a filler.
\end{enumerate}
In both cases the filler is computed by the same algorithm: choose any
covering fiber, take its local section, and verify uniqueness via the
compatibility data. \hfill\(\square\)
\end{proof}

\begin{corollary}
The compiler rule \texttt{\_compile\_cubical\_filling} handles both
\textsf{RefinementDescent} and \textsf{HComp} proof terms.  The
only difference is whether the face predicates are Z3-decidable
refinement predicates (descent) or interval names (hcomp).
\end{corollary}


% ─────────────────────────────────────────────────────────
\section{The Two Dimension Systems}
\label{sec:two-dimensions}

\begin{definition}[The Two Dimension Systems of \(\mathrm{C}^4\)]
\label{def:two-dimensions}
\(\mathrm{C}^4\) has two independent dimension systems, each contributing
a class of ``directions'' in which proofs can vary:

\medskip
\begin{center}
\begin{tabular}{lll}
\toprule
& \textbf{Cubical dimensions} & \textbf{Cohomological dimensions} \\
\midrule
Parameterized by & Interval variables \(i : \interval\) &
  Refinement predicates \(\varphi : \alpha \to \mathsf{Prop}\) \\
0-cells & Programs & Input values \\
1-cells & Equality proofs (paths) & Fiber restrictions \\
2-cells & Proof coherence (squares) & Cover overlaps \\
Composition & \(\hfilling\) (Kan) & Descent (gluing) \\
Transport & Along type paths & Along predicate implications \\
Obstruction & Kan condition & \(\check{H}^1 = 0\) \\
Decidability & Structural & Z3 \\
Algebra & \(\interval\): De Morgan lattice &
  \(\mathcal{P}(\alpha \to \mathsf{Prop})\): Boolean algebra \\
\bottomrule
\end{tabular}
\end{center}
\end{definition}

A proof over a function with \(m\) proof steps and \(n\) branches
lives in \(\interval^m \times \site^n\): the cubical dimensions
parameterize \emph{how} we prove; the cohomological dimensions
parameterize \emph{for which inputs} we prove it.

\begin{theorem}[Algebraic Synergy of the Two Dimensions]
\label{thm:algebraic-synergy}
The cubical interval connections \(\wedge_i, \vee_i\) on \(\interval\)
and the cover operations \(\wedge, \vee\) on \(\mathcal{P}(\alpha \to \mathsf{Prop})\)
satisfy the same De Morgan laws:
\[
  \neg(i \wedge_i j) = \neg_i i \vee_i \neg_i j,
  \qquad
  \neg(\varphi \wedge \psi) = \neg\varphi \vee \neg\psi.
\]
Consequently, all cover simplifications---absorption, distribution,
idempotence, complementation---are valid in both dimensions
simultaneously.
\end{theorem}

\begin{proof}
The cubical De Morgan laws hold in the model of cubical type theory
by the face map identities.  The predicate De Morgan laws hold
propositionally in any Boolean algebra.  Both algebras are free
Boolean algebras (the interval on one generator, the predicate algebra
on the set of atomic predicates), and their quotient by the respective
coherence conditions is the same De Morgan algebra.
\hfill\(\square\)
\end{proof}


% ─────────────────────────────────────────────────────────
\section{Design Principles}
\label{sec:design-principles}

The following principles follow from the fundamental identification
and the two-dimension structure:

\begin{enumerate}
\item \textbf{Refinement site.}  The site objects are
  Z3-decidable refinement predicates \(\varphi : \alpha \to \mathsf{Prop}\),
  not syntactic type tags.  Site morphisms are Z3-provable implications
  \(\varphi \Rightarrow \psi\).  Every program-level type annotation
  is simultaneously a type-theoretic refinement (typing judgment)
  and a cohomological face (proof decomposition point).

\item \textbf{Descent \(=\) HComp.}
  There is one filling rule, not two.  It is parameterized over
  whether the face predicates are interval names or refinement
  predicates.

\item \textbf{Strict superset of CIC.}  Every CIC judgment is valid
  in \(\mathrm{C}^4\) (trivial one-fiber site).  \(\mathrm{C}^4\) is a
  conservative extension.

\item \textbf{Cubical paths are native.}  Path types \(\Path_A\, a\, b\)
  are defined via the interval, not inductive identity types.
  Univalence and function extensionality hold definitionally.

\item \textbf{Hedberg collapse.}  For Python verification, all types
  have decidable equality (integers, strings, booleans, and their
  combinations).  By Hedberg's theorem, \(\check{H}^1 = 0\) for all
  propositional coefficients.  Thus descent is \emph{always} applicable:
  there is no obstruction to gluing.

\item \textbf{Trust as a sheaf.}  Each fiber carries a trust provenance
  set.  The global proof's trust is the union over fibers.

\item \textbf{OTerms are first-class.}  The Python denotational
  semantics inhabit a universe \(\mathcal{U}_\PyObj\).
  CCTT axioms (D1--D24 etc.) are equalities in that universe.
\end{enumerate}


% ─────────────────────────────────────────────────────────
\section{Syntax}
\label{sec:syntax}

\subsection{The Refinement Site}

\begin{definition}[Refinement Site]
\label{def:refinement-site}
Fix a base type \(\alpha : \mathcal{U}_0\).  The
\textbf{refinement site} \(\site_\alpha\) has:
\begin{itemize}
\item \textbf{Objects:}  Z3-decidable predicates
  \(\varphi : \alpha \to \mathsf{Prop}\) with a decidability witness
  \(\delta_\varphi : \forall x.\, \varphi(x) \vee \neg\varphi(x)\).
\item \textbf{Morphisms:}
  \((\varphi \to_{\site} \psi) \;\coloneqq\;
  \forall x.\, \varphi(x) \Rightarrow \psi(x)\),
  the set of Z3-provable implications.
\item \textbf{Grothendieck topology:}  A \textbf{cover} of the
  terminal object is a finite family
  \(\covers = \{\varphi_1, \ldots, \varphi_n\}\) with
  \(\forall x.\,\bigvee_i \varphi_i(x)\).
\end{itemize}
The site satisfies the axioms of a Grothendieck topology:
stability under base change and local character.
\end{definition}

\begin{remark}
When the base type is \(\mathbb{Z}\), the site objects are integer
predicates such as \(\{x > 0\}\), \(\{x \geq 0\}\),
\(\{\text{even}(x)\}\).  The morphisms are implications provable
by linear arithmetic.  The covers correspond to case splits in Python
(\texttt{if/elif/else}).
\end{remark}

\subsection{Sorts and Universes}

\(\mathrm{C}^4\) has a cumulative hierarchy of universes, each
\emph{site-decorated}:
\[
  \mathsf{Prop} \;\subset\; \mathcal{U}_0 \;\subset\;
  \mathcal{U}_1 \;\subset\; \mathcal{U}_2 \;\subset\; \cdots
\]

\begin{definition}[Site-Decorated Universe]
A \textbf{site-decorated universe} \(\mathcal{U}_i^{\site}\) is a
universe of sheaves over the site \(\site\): its types are functors
\(A : \site^{\mathrm{op}} \to \mathcal{U}_i\), so that for each object
\(\varphi \in \site\), we have \(A(\varphi) : \mathcal{U}_i\), and
for each morphism \(f : \psi \to \varphi\), there is a restriction
map \(\restr{}{f} : A(\varphi) \to A(\psi)\).
\end{definition}

The OTerm universe \(\mathcal{U}_\PyObj\) is site-decorated over
the refinement site for Python programs.

\subsection{Term Grammar}

\[
\begin{array}{rcll}
  t, A, B & ::= & x & \text{variable} \\
  & \mid & \mathcal{U}_i^{\site} & \text{universe at level \(i\) over site \(\site\)} \\
  & \mid & \Pi(x : A).\, B & \text{dependent product} \\
  & \mid & \lambda(x : A).\, t \mid t\; u & \text{abstraction / application} \\
  & \mid & \Sigma(x : A).\, B \mid (t, u) \mid \pi_1 t \mid \pi_2 t
    & \text{dependent sum} \\
  & \mid & \mathsf{Ind}(\Delta) \mid \mathsf{rec}_\Delta(t, \vec{c})
    & \text{inductive types} \\[4pt]
  & \mid & \interval \mid \mathsf{0}_\interval \mid \mathsf{1}_\interval
    & \text{interval (cubical)} \\
  & \mid & \Path_A\, a\, b \mid \langle i \rangle\, t \mid t\; r
    & \text{path type, intro, elim} \\
  & \mid & \transp^i\, A\, \varphi\, t
    & \text{transport (cubical)} \\
  & \mid & \hfilling^i\, A\, [\vec{f} \mapsto s_f]\, t_0
    & \text{homogeneous composition} \\
  & \mid & \Glue[\varphi \mapsto (T, e)]\, A
    & \text{glue type} \\[4pt]
  & \mid & \{x : A \mid \varphi(x)\}
    & \text{refinement type (cohomological)} \\
  & \mid & \langle a, h \rangle_\varphi \mid r.1 \mid r.2
    & \text{refinement intro / elim} \\
  & \mid & \restr{t}{\varphi}
    & \text{fiber restriction} \\
  & \mid & \mathsf{desc}(\covers, \{t_\alpha\}, \{p_{\alpha\beta}\})
    & \text{descent / global section} \\[4pt]
  & \mid & \mathsf{import}(\mathit{name})
    & \text{Mathlib axiom import} \\
  & \mid & \den{P}
    & \text{OTerm embedding} \\
\end{array}
\]

\subsection{Contexts}

A context \(\Gamma\) is a list of typed variable bindings:
\[
  \Gamma \;::=\; \cdot \;\mid\; \Gamma,\, x : A \;\mid\;
  \Gamma,\, x :_\varphi A
\]
where \(x :_\varphi A\) binds \(x : A\) under the fiber hypothesis
\(\varphi(x)\) --- formally: \(x : \{y : A \mid \varphi(y)\}\).

The fiber hypothesis binding is the key to \textbf{hypothesis pushing}
(Synergy~6 of the atlas): within a fiber proof, the predicate \(\varphi\)
is freely available as an assumption.


% ─────────────────────────────────────────────────────────
\section{Judgments}
\label{sec:judgments}

\(\mathrm{C}^4\) has six judgment forms:

\begin{enumerate}
\item \(\Gamma \vdash A : \mathcal{U}_i^{\site}\)
  \hfill (``\(A\) is a sheaf-type at level \(i\)'')
\item \(\Gamma \vdash t : A\)
  \hfill (``\(t\) is a global section of \(A\)'')
\item \(\Gamma \vdash A \equiv B : \mathcal{U}_i^{\site}\)
  \hfill (definitional type equality)
\item \(\Gamma \vdash t \equiv u : A\)
  \hfill (definitional term equality)
\item \(\Gamma \vdash_\varphi t : A\)
  \hfill (``\(t\) is a local section of \(A\) on fiber \(\varphi\)'')
\item \(\Gamma \vdash t :_\mathcal{T} A\)
  \hfill (``\(t\) has type \(A\) with trust provenance \(\mathcal{T}\)'')
\end{enumerate}

Judgments (1)--(4) are standard CIC-like.  Judgment~(5) is the
cohomological fiber judgment: it allows reasoning locally on fiber
\(\varphi\) with \(\varphi\) available as a hypothesis.
Judgment~(6) tracks the trust provenance of every typing derivation.


% ─────────────────────────────────────────────────────────
\section{Typing Rules}
\label{sec:typing-rules}

\subsection{Standard CIC Rules}

The standard CIC rules are unchanged.  Key rules for reference:

\begin{gather}
\frac{}{\Gamma \vdash \mathcal{U}_i : \mathcal{U}_{i+1}}
\quad (\mathcal{U}\text{-form})
\tag{Univ}
\\[6pt]
\frac{\Gamma \vdash A : \mathcal{U}_i \qquad
      \Gamma, x : A \vdash B : \mathcal{U}_j}
     {\Gamma \vdash \Pi(x:A).\, B : \mathcal{U}_{\max(i,j)}}
\tag{$\Pi$F}
\\[6pt]
\frac{\Gamma, x : A \vdash t : B}
     {\Gamma \vdash \lambda(x:A).\, t : \Pi(x:A).\, B}
\tag{$\Pi$I}
\\[6pt]
\frac{\Gamma \vdash f : \Pi(x:A).\, B \qquad \Gamma \vdash a : A}
     {\Gamma \vdash f\; a : B[x := a]}
\tag{$\Pi$E}
\end{gather}

\subsection{Cubical Rules}

\begin{gather}
\frac{\Gamma \vdash A : \mathcal{U}_i \qquad
      \Gamma \vdash a : A \qquad \Gamma \vdash b : A}
     {\Gamma \vdash \Path_A\, a\, b : \mathcal{U}_i}
\tag{PathF}
\\[6pt]
\frac{\Gamma, i : \interval \vdash t : A \qquad
      t[i := 0] \equiv a \qquad t[i := 1] \equiv b}
     {\Gamma \vdash \langle i \rangle\, t : \Path_A\, a\, b}
\tag{PathI}
\\[6pt]
\frac{\Gamma \vdash p : \Path_A\, a\, b \qquad
      \Gamma \vdash r : \interval}
     {\Gamma \vdash p\; r : A \qquad p\; 0 \equiv a \qquad p\; 1 \equiv b}
\tag{PathE}
\\[6pt]
\frac{\Gamma, i : \interval \vdash A : \mathcal{U}_i \qquad
      \Gamma \vdash t : A[i := 0]}
     {\Gamma \vdash \transp^i\, A\, \mathsf{1}_\varphi\, t : A[i := 1]}
\tag{Transp}
\\[6pt]
\frac{
  \Gamma, i : \interval \vdash A : \mathcal{U}_i
  \qquad
  \Gamma \vdash t_0 : A[i := 0]
  \qquad
  \forall k.\;
  \Gamma, i : \interval, \psi_k \vdash u_k : A
  \quad
  u_k[i := 0] \equiv t_0
}
{
  \Gamma \vdash \hfilling^i\, A\, [\psi_k \mapsto u_k]\, t_0
  : A[i := 1]
}
\tag{HComp}
\end{gather}

\subsection{Refinement Type Rules}

These are new relative to standard CIC:

\begin{gather}
\frac{\Gamma \vdash A : \mathcal{U}_i \qquad
      \Gamma, x : A \vdash \varphi(x) : \mathsf{Prop} \qquad
      \text{Z3 can decide } \varphi}
     {\Gamma \vdash \{x : A \mid \varphi(x)\} : \mathcal{U}_i}
\tag{RefinF}
\\[6pt]
\frac{\Gamma \vdash a : A \qquad \Gamma \vdash h : \varphi(a)}
     {\Gamma \vdash \langle a, h \rangle_\varphi :
      \{x : A \mid \varphi(x)\}}
\tag{RefinI}
\\[6pt]
\frac{\Gamma \vdash r : \{x : A \mid \varphi(x)\}}
     {\Gamma \vdash r.1 : A \qquad \Gamma \vdash r.2 : \varphi(r.1)}
\tag{RefinE}
\end{gather}

\begin{remark}[Refinement Types are Cubical Faces]
A refinement type \(\{x : A \mid \varphi(x)\}\) is simultaneously:
\begin{enumerate}
\item A subtype in the type-theoretic sense (elements satisfying \(\varphi\)).
\item A cubical face in the proof cube (restricting the proof to inputs
  satisfying \(\varphi\)).
\item A sheaf restriction \(\restr{\mathcal{F}}{\{\varphi\}}\) for local
  verification.
\end{enumerate}
This triple role (Synergy~1 of the atlas) is why adding a type annotation
automatically provides both a proof decomposition point and a Z3 hypothesis.
\end{remark}

\subsection{The Unified Filling Rule}

The descent rule and the HComp rule are unified as:

\begin{definition}[The Unified Filling Rule]
\label{def:unified-filling}
Given:
\begin{itemize}
\item A type \(A\) (sheaf over the site),
\item A cover \(\covers = \{\varphi_1, \ldots, \varphi_n\}\)
  (either interval faces or refinement predicates),
\item Local sections \(t_\alpha : \restr{A}{\varphi_\alpha}\)
  (either face terms or fiber proofs),
\item Compatibility proofs \(p_{\alpha\beta} :
  \restr{t_\alpha}{\varphi_\alpha \wedge \varphi_\beta}
  \equiv
  \restr{t_\beta}{\varphi_\alpha \wedge \varphi_\beta}\),
\item Exhaustiveness \(\vdash \bigvee_\alpha \varphi_\alpha\) (or
  the cover is the full boundary in the cubical case):
\end{itemize}
\[
\frac{
  \begin{array}{c}
    \covers = \{\varphi_\alpha\}_{\alpha \in I}
    \qquad
    \forall x.\, \bigvee_\alpha \varphi_\alpha(x)
    \\[4pt]
    \forall \alpha.\;
    \Gamma \vdash_{\varphi_\alpha} t_\alpha : \restr{A}{\varphi_\alpha}
    \\[4pt]
    \forall \alpha, \beta.\;
    \Gamma \vdash p_{\alpha\beta}
    : \restr{t_\alpha}{\varphi_\alpha \wedge \varphi_\beta}
    \equiv
    \restr{t_\beta}{\varphi_\alpha \wedge \varphi_\beta}
    \\[4pt]
    \check{H}^1(\covers, \underline{\mathrm{Aut}}(A)) = 0
  \end{array}
}
{
  \Gamma \vdash
  \mathsf{fill}(\covers, \{t_\alpha\}, \{p_{\alpha\beta}\}) : A
}
\tag{Fill}
\]
\end{definition}

When \(\covers\) consists of interval faces \([i_k = \varepsilon]\),
this instantiates to \(\hfilling\).  When \(\covers\) consists of
refinement predicates, this instantiates to \(\mathsf{desc}\).

\begin{theorem}[Hedberg Collapse]
\label{thm:hedberg-collapse}
For all sheaves \(A\) with propositional fibers (i.e.,
\(\forall \varphi.\; A(\varphi)\) is a proposition),
\[
  \check{H}^1(\covers, \underline{\mathrm{Aut}}(A)) = 0
\]
for every cover \(\covers\).  Therefore the \(\mathrm{C}^4\)
filling rule \textup{(Fill)} is applicable without an explicit
cohomological obstruction check.
\end{theorem}

\begin{proof}
Hedberg's theorem (see Chapter~\ref{ch:c4-metatheory},
Theorem~\ref{thm:hedberg-main}) states that types with decidable equality
are sets (proof-irrelevant for equalities).  The sheaf
\(\underline{\mathrm{Aut}}(A)\) has values \(\mathrm{Aut}(A(\varphi))\),
the automorphism group of the local fiber.  For propositional fibers,
the only automorphism is the identity, so
\(\underline{\mathrm{Aut}}(A)\) is the trivial sheaf.
The first \(\check{C}\)ech cohomology of a trivial sheaf is zero.
\hfill\(\square\)
\end{proof}

\begin{corollary}
In the C\(^4\) verification setting for Python (all types have
Z3-decidable equality), the cohomological obstruction check is never
needed.  The filling rule fires unconditionally.
\end{corollary}

\subsection{Restriction Rules}

\begin{gather}
\frac{\Gamma \vdash t : A \qquad \varphi \in \site}
     {\Gamma \vdash_\varphi \restr{t}{\varphi} :
      \restr{A}{\varphi}}
\tag{Res}
\\[6pt]
\frac{\Gamma \vdash_\varphi t : A \qquad
      \Gamma \vdash (\psi \Rightarrow \varphi) : \mathsf{Prop}}
     {\Gamma \vdash_\psi \restr{t}{(\psi \hookrightarrow \varphi)} :
      \restr{A}{\psi}}
\tag{Pull}
\end{gather}

Rule~(Res) restricts any global section to a fiber.
Rule~(Pull) is transport along a site morphism (an implication):
it corresponds to the cubical transport rule where the ``path'' is
the implication \(\psi \Rightarrow \varphi\).

\subsection{Trust Typing Rule}

\begin{gather}
\frac{\Gamma \vdash t : A \qquad \mathcal{T} = \mathrm{trust}(t)}
     {\Gamma \vdash t :_{\mathcal{T}} A}
\tag{Trust}
\\[6pt]
\frac{\Gamma \vdash t :_{\mathcal{T}_1} A \qquad
      \Gamma \vdash u :_{\mathcal{T}_2} B[x := t]}
     {\Gamma \vdash (t, u) :_{\mathcal{T}_1 \cup \mathcal{T}_2}
      \Sigma(x : A).\, B}
\tag{Trust$\Sigma$}
\end{gather}

The trust set is the upward closure of the trust sources used in
the derivation.  This gives trust its sheaf structure: the trust of
a descent proof is the union of the trust of each fiber proof.

\subsection{F*-Style Binding Rule}

The F*-style annotation compatibility rule (Synergy~34 of the atlas)
is a structural rule that checks annotation--code agreement:

\begin{gather}
\frac{
  \begin{array}{c}
    \Gamma \vdash_\mathrm{ann} P : \Pi(x : A).\, \mathsf{Prop}
    \quad \text{(annotation, a predicate on the output)}
    \\[3pt]
    \Gamma \vdash_\mathrm{code} f : \Pi(x : A).\, B
    \quad \text{(code, a function)}
    \\[3pt]
    \forall \varphi \in \covers.\;
    \Gamma \vdash_\varphi P(\restr{f}{\varphi}\, x) : \mathsf{Prop}
    \quad \text{(annotation matches code per fiber)}
    \\[3pt]
    \covers \text{ exhaustive over } A
  \end{array}
}
{
  \Gamma \vdash \mathsf{check\_binding}(P, f, \covers) :
  \Pi(x : A).\, P(f\, x)
}
\tag{Bind}
\end{gather}

This rule expresses that an annotation is only accepted when it can
be witnessed per-fiber (cohomological) by proofs that compose into a
global derivation (cubical path composition).


% ─────────────────────────────────────────────────────────
\section{Computation Rules}
\label{sec:computation-rules}

\subsection{Standard \(\beta\)-Reductions}

\begin{align}
  (\lambda(x : A).\, t)\; u
  &\;\rightsquigarrow\; t[x := u]
  \tag{$\beta$}
  \\
  \pi_1(a, b) &\;\rightsquigarrow\; a
  \tag{$\Sigma\beta_1$}
  \\
  \pi_2(a, b) &\;\rightsquigarrow\; b
  \tag{$\Sigma\beta_2$}
\end{align}

\subsection{Cubical \(\beta\)-Reductions}

\begin{align}
  (\langle i \rangle\, t)\; r
  &\;\rightsquigarrow\; t[i := r]
  \tag{Path-$\beta$}
  \\
  \transp^i\, A\, \mathsf{1}\, t
  &\;\rightsquigarrow\; t
  \qquad\text{(transport on constant family)}
  \tag{Transp-const}
\end{align}

\subsection{Restriction Distribution}

Restriction is a left-exact functor; it distributes over all
type formers:
\begin{align}
  \restr{(\lambda x.\, t)}{\varphi}
  &\;\rightsquigarrow\; \lambda x.\, \restr{t}{\varphi}
  \tag{Res-$\lambda$}
  \\
  \restr{(t, u)}{\varphi}
  &\;\rightsquigarrow\; (\restr{t}{\varphi}, \restr{u}{\varphi})
  \tag{Res-$\Sigma$}
  \\
  \restr{(\langle i \rangle\, t)}{\varphi}
  &\;\rightsquigarrow\; \langle i \rangle\, \restr{t}{\varphi}
  \tag{Res-Path}
  \\
  \restr{(\transp^i\, A\, h\, t)}{\varphi}
  &\;\rightsquigarrow\;
  \transp^i\, (\restr{A}{\varphi})\, \restr{h}{\varphi}\,
  \restr{t}{\varphi}
  \tag{Res-Transp}
\end{align}

\subsection{Descent \(\beta\)-Reduction (Key Rule)}

\begin{equation}
  \restr{\mathsf{fill}(\covers, \{t_\alpha\}, \{p_{\alpha\beta}\})}{\varphi_k}
  \;\rightsquigarrow\; t_k
  \tag{Fill-$\beta$}
\end{equation}

This rule asserts: descending to a global section and then restricting
to a fiber gives back the local section.  This is the descent
\(\beta\)-reduction, and it is the dual of the
Path-\(\beta\) rule.

\begin{remark}
The reduction system is terminating because (a) restriction
pushes inward (strictly decreasing the depth of outermost
restriction operators), (b) Fill-\(\beta\) eliminates
\(\mathsf{fill}\) nodes, and (c) the CIC and cubical fragments
terminate independently.  See Chapter~\ref{ch:c4-metatheory}
for the formal argument via a lexicographic measure.
\end{remark}


% ─────────────────────────────────────────────────────────
\section{The Cover Algebra}
\label{sec:cover-algebra}

The refinement site has a rich algebra of covers.  These correspond
to operations on cubical face systems (Synergy~10 of the atlas):

\begin{definition}[Cover Operations]
Let \(\covers = \{\varphi_i\}\) and \(\mathcal{V} = \{\psi_j\}\)
be covers of \(\alpha\).
\begin{align*}
  \text{Refinement:}\quad
  &\covers' \leq \covers
  \;\Leftrightarrow\;
  \forall j.\; \exists i.\; \varphi'_j \Rightarrow \varphi_i \\
  \text{Product:}\quad
  &\covers \times \mathcal{V}
  \;\coloneqq\; \{\varphi_i \wedge \psi_j\}_{i,j} \\
  \text{Join:}\quad
  &\covers \vee \mathcal{V}
  \;\coloneqq\; \{\varphi_i\}_i \cup \{\psi_j\}_j \\
  \text{Restriction:}\quad
  &\restr{\covers}{\psi}
  \;\coloneqq\; \{\varphi_i \wedge \psi\}_i
\end{align*}
\end{definition}

\begin{proposition}[Cover algebra is a locale]
The cover operations satisfy the axioms of a locale (complete
Heyting algebra), with \(\leq\) as the partial order.
The product \(\times\) is the meet, the join \(\vee\) is the join,
and restriction \(\restr{\cdot}{\psi}\) is the right adjoint
of product with \(\{\psi\}\).
\end{proposition}

\begin{proof}
All operations are Z3-computable (predicates with decidable logic).
The locale axioms reduce to standard propositional identities:
distribution, modularity, and monotonicity, all of which hold
in a Boolean algebra.  \hfill\(\square\)
\end{proof}

\textbf{Cubical interpretation.}
Each cover operation corresponds to a cubical face operation:
product \(\covers \times \mathcal{V}\) corresponds to the product
cube \(\interval^m \times \interval^n = \interval^{m+n}\);
restriction \(\restr{\covers}{\psi}\) corresponds to the face
evaluation \(\partial^\varepsilon_k\); the join \(\covers \vee \mathcal{V}\)
corresponds to adding cubical dimensions.

This is why the algebraic synergy of Theorem~\ref{thm:algebraic-synergy}
is not superficial: the cover algebra and the face algebra are both
quotients of the same free Boolean algebra.


% ─────────────────────────────────────────────────────────
\section{The OTerm Universe}
\label{sec:oterm-universe}

\begin{definition}[OTerm Universe]
The \textbf{OTerm universe} \(\mathcal{U}_\PyObj\) is a
site-decorated universe in \(\mathrm{C}^4\) whose:
\begin{enumerate}
\item \textbf{Types} are OTerms (elements of the free algebra
  \(\PyObj\) generated by \(\mathsf{OVar}, \mathsf{OLit}, \mathsf{OOp},
  \mathsf{OFold}, \ldots\)).
\item \textbf{Site} is the refinement site for Python inputs:
  objects are Z3-decidable predicates over \(\mathsf{Value}\),
  morphisms are Z3-provable implications.
\item \textbf{Axiom equalities}: each CCTT axiom
  \(D_k\) or \(P_k\) gives a path in \(\mathcal{U}_\PyObj\).
\end{enumerate}
\end{definition}

The OTerm universe unifies the two dimension systems for Python programs:
\begin{itemize}
\item Cubical paths in \(\mathcal{U}_\PyObj\) are program equivalences
  (chains of axiom applications).
\item Cohomological fibers in \(\mathcal{U}_\PyObj\) are type-class
  branches (\texttt{isinstance} cases).
\end{itemize}

A proof that two Python functions are equivalent is a path
\(\Path_{\mathcal{U}_\PyObj}\, \den{f}\, \den{g}\) in the OTerm
universe, which decomposes via descent into per-fiber paths
(per-branch equivalences) satisfying the overlap conditions.


% ─────────────────────────────────────────────────────────
\section{Inductive Types and Higher Inductive Types}
\label{sec:inductive-types}

$\mathrm{C}^4$ supports all CIC inductive types and additionally
supports \textbf{higher inductive types} (HITs) with path
constructors.  The \textbf{PyComp HIT} models the denotational
semantics of Python:

\begin{example}[PyComp as a HIT]
\vspace{-4pt}
\begin{align*}
  \mathsf{data}\; \PyObj &: \mathcal{U}_0 \;\mathsf{where} \\
  \mathsf{OVar}(s)   &: \PyObj \quad (s : \mathsf{String}) \\
  \mathsf{OLit}(v)   &: \PyObj \quad (v : \mathsf{Value}) \\
  \mathsf{OOp}(f, \vec{a}) &: \PyObj \quad \text{(built-in operator)} \\
  \mathsf{OFold}(\oplus, z, xs) &: \PyObj \quad \text{(fold)} \\
  \mathsf{OCase}(c, t, e) &: \PyObj \quad \text{(conditional)} \\
  \mathsf{OFix}(x, b) &: \PyObj \quad \text{(fixed point)} \\
  &\vdots \\
  \mathsf{D1}     &: \Path_\PyObj\,
    (\mathsf{OOp}(+, [a, b]))\;
    (\mathsf{OOp}(+, [b, a]))
    \quad \text{(commutativity of \texttt{+})}
\end{align*}

The 24 dichotomy path constructors D1--D24 are equalities that
hold definitionally in $\mathrm{C}^4$.  Any OTerm equality derivable
from these paths via path composition, transport, and descent is
automatically provable.
\end{example}

The HITs interact with the covering structure: the fibration of
\(\PyObj\) over the refinement site (predicates on Python values)
gives each HIT path constructor a fiber structure.  A path
constructor that only applies on fiber \(\varphi\) (e.g., D1 applies
on \(\{\mathsf{isinstance}(a, \mathsf{int})\}\)) is a
\emph{local} path in the OTerm universe.


% ─────────────────────────────────────────────────────────
\section{Relationship to CIC, Cubical TT, and F*}
\label{sec:relationships}

\subsection{CIC}

\begin{theorem}[Conservative Extension over CIC]
\label{thm:conservativity-cic}
$\mathrm{C}^4$ is a conservative extension of CIC: a judgment is
derivable in CIC if and only if it is derivable in $\mathrm{C}^4$
via the trivial one-fiber site.
\end{theorem}

\begin{proof}[Proof sketch]
Forward: CIC rules are a subset of $\mathrm{C}^4$ rules (the trivial
site \(\{\mathrm{True}\}\) makes all restriction and descent rules
degenerate).

Reverse: the $\mathrm{C}^4 \to \mathrm{CIC}$ erasure functor
removes fiber annotations and replaces descent with the identity.
It preserves typing because the erasure of a $\mathrm{C}^4$
derivation is a valid CIC derivation.  The novel rules (restriction,
descent, refinement) all erase to the standard CIC rules for the
trivial cover.  \hfill\(\square\)
\end{proof}

\subsection{Cubical Type Theory}

\begin{theorem}[Conservative Extension over Cubical TT]
$\mathrm{C}^4$ is a conservative extension of the cubical type theory
of Cohen--Coquand--Huber--M\"ortberg.  In particular, univalence and
function extensionality hold as computed equalities, not axioms.
\end{theorem}

\begin{proof}[Proof sketch]
The cubical fragment of $\mathrm{C}^4$ is exactly CCHM cubical type
theory, with one addition: the face predicates \(\psi_k\) in the HComp
rule may now be \emph{either} interval names \emph{or} refinement
predicates.  The CCHM rules are recovered by taking only interval
names.  The addition of refinement predicates does not change the
behavior of the pure interval fragment.  \hfill\(\square\)
\end{proof}

\subsection{F*}

F* uses a refinement-type system with Z3-based discharge.  Compared
to $\mathrm{C}^4$:
\begin{enumerate}
\item F* does not have the cubical interval or path types.  Equality
  is propositional (\(\mathsf{Prop}\)-valued), not proof-relevant.
\item F* does not have the sheaf structure.  Refinement types in F*
  do not have a cover algebra or descent rule.
\item F* does not have the Hedberg collapse metatheorem --- it must
  be proved manually for each type where it is needed.
\item The F*-style binding rule (\textsc{Bind}) is captured in $\mathrm{C}^4$
  as a structural rule (not a metalevel annotation language), making it
  first-class in the type theory.
\end{enumerate}
$\mathrm{C}^4$ strictly extends F*'s expressiveness by adding (a) path
types for code equivalence proofs, (b) the cover algebra for
interprocedural composition, and (c) the unified descent=HComp rule
for efficient proof assembly.
